\documentclass[../../main.tex]{subfiles}
\graphicspath{{\subfix{../../image/}}}

\begin{document}

\section{几何概率}\label{section:ch01-几何概率}

\subsection*{一、例子与计算公式}

在古典概型中利用等可能性的概念，成功地计算了某一类问题的概率；不过，古典概型要求可能场合的总数必须有限，因此历史上有不少人企图把这种做法推广到有无限多结果而又有某种等可能性的场合。这类问题一般可以通过几何方法来求解。

先从几个简单的例子开始。

\begin{example}[等待广播]\label{example:ch01-等待广播}
某人午觉醒来，发觉表停了，他打开收音机，想听电台整点报时，求他等待的时间短于 $10$ 分钟的概率。
\end{example}

\begin{example}[海域石油]\label{example:ch01-海域石油}
如果在一个 $5\times 10^4\,\mathrm{km}^2$ 的海域里有表面积达 $40\,\mathrm{km}^2$ 的大陆架贮藏着石油，假如在这海域里随意选定一点钻探，问钻到石油的概率是多少？
\end{example}

\begin{example}[大肠杆菌]\label{example:ch01-大肠杆菌}
在 $400\,\mathrm{mL}$ 自来水中有一个大肠杆菌，今从中随机取出 $2\,\mathrm{mL}$ 水样放到显微镜下观察，求发现大肠杆菌的概率。
\end{example}

一种相当自然的答案是认为例~\ref{example:ch01-等待广播} 所求的概率等于 $\frac{1}{6}$，例~\ref{example:ch01-海域石油} 中钻到石油的概率等于 $\frac{8}{10\,000}$，而例~\ref{example:ch01-大肠杆菌} 所求的概率等于 $\frac{1}{200}$。在求这些概率时，我们事实上是利用了几何的方法，并假定了某种等可能性。

在例~\ref{example:ch01-等待广播} 中，因为电台每小时报时一次，我们自然认为这个人打开收音机时处于两次报时之间，例如 $(13:00,14:00)$，而且取各点的可能性一样，要遇到等待时间短于 $10$ 分钟，只有当他打开收音机的时间正好处于 $13:50$ 至 $14:00$ 之间才有可能，相应的概率是 $\frac{10}{60}=\frac{1}{6}$。

在例~\ref{example:ch01-海域石油} 中，由于选点的随机性，可以认为该海域中各点被选中的可能性是一样的，因而所求概率自然认为等于贮油海域的面积与整个海域面积之比，即等于 $\frac{40}{50\,000}$。

同样地，例~\ref{example:ch01-大肠杆菌} 中由于取水样的随机性，所求概率等于水样的体积与总体积之比 $\frac{2}{400}$。

总之，在这类问题中，试验的可能结果是某区域 $\Omega$ 中的一个点。这个区域可以是一维的，也可以是二维的，还可以是三维的，甚至可以是 $n$ 维的，这时不管是可能结果全体或是我们所感兴趣的结果都是无限的。因而等可能性是通过下列方式来赋予意义的：落在某区域 $g$ 的概率与区域 $g$ 的测度（长度、面积、体积等等）成正比并且与其位置及形状无关。

因此，若以 $A_g$ 记“在区域 $\Omega$ 中随机地取一点，而该点落在区域 $g$ 中”这一事件，则其概率定义为
\begin{align}
P(A_g)=\frac{g\text{ 的测度}}{\Omega\text{ 的测度}}\label{align:ch01-几何概率定义}
\end{align}
据此定义，则上述诸例之解是明显的。下面再举一个例子。

\begin{example}[会面问题]\label{example:ch01-会面问题}
两人相约 $7$ 点到 $8$ 点在某地会面，先到者等候另一人 $20$ 分钟，过时就可离去，试求这两人能会面的概率。

\textbf{解}~~以 $x,y$ 分别表示两人到达时刻，则会面的充要条件为
\begin{align*}
|x-y|\leqslant 20
\end{align*}
这是一个几何概率问题，可能的结果全体是边长为 $60$ 的正方形里的点，能会面的点的区域用阴影标出（图~\ref{figure:ch01-会面问题}）。所求概率为
\begin{align*}
p=\frac{60^2-40^2}{60^2}=\frac{5}{9}
\end{align*}
\end{example}

\begin{figure}[H]
\centering
\includegraphics[width=0.5\textwidth]{ch01-meeting.png}
\caption{会面问题}\label{figure:ch01-会面问题}
\end{figure}

\subsection*{二、Buffon 问题}

$1777$ 年法国科学家 Buffon（$1701$—$1788$）提出了下列著名问题，这是几何概率的一个早期例子。

\begin{example}[投针问题]\label{example:ch01-投针}
平面上画着一些平行线，它们之间的距离都等于 $a$，向此平面任投一长度为 $l$（$l<a$）的针，试求此针与任一平行线相交的概率。

\textbf{解}~~以 $x$ 表示针的中点到最近的一条平行线的距离，$\varphi$ 表示针与平行线的交角。针与平行线的位置关系见图~\ref{figure:ch01-buffon}。

\begin{figure}[H]
\centering
\includegraphics[width=0.8\textwidth]{ch01-buffon.png}
\caption{Buffon 问题}\label{figure:ch01-buffon}
\end{figure}

显然有 $0\leqslant x\leqslant\frac{a}{2}$，$0\leqslant\varphi\leqslant\pi$，以 $G$ 表示边长为 $\frac{a}{2}$ 及 $\pi$ 的长方形。为使针与平行线相交，必须 $x\leqslant\frac{l}{2}\sin\varphi$，满足这个关系式的区域记为 $g$，在图~\ref{figure:ch01-buffon} 中用阴影表出，所求的概率为
\begin{align}
p=\frac{g\text{ 的面积}}{G\text{ 的面积}}=\frac{\frac{1}{2}\int_0^\pi l\sin\varphi\,\mathrm{d}\varphi}{\frac{1}{2}a\pi}=\frac{2l}{\pi a}\label{align:ch01-buffon}
\end{align}

由于最后的答案与 $\pi$ 有关，因此 Buffon 设想利用它来计算 $\pi$ 的数值，其方法是投针 $N$ 次，计算针与线相交的次数 $n$，再以频率值 $\frac{n}{N}$ 作为概率 $p$ 之值代入~\eqref{align:ch01-buffon}，求得
\begin{align}
\pi=\frac{2lN}{an}\label{align:ch01-buffon-pi}
\end{align}

粗看这是一个笨办法，耗力费时，而且很难达到当时用数学方法已算得小数点后一百多位精确数值这样一个精度。但仔细考量则会发现这是一个了不起的创意，它提出了一个全新的计算方案：建立一个概率模型，它与某些我们感兴趣的量——这里是常数 $\pi$——有关，然后设计适当的随机试验，并通过这个试验的结果来确定这些量。

据说 Buffon 曾亲自做过试验，可惜结果没有留传下来。

历代欣赏 Buffon 提议的大有人在，有些还真正做了试验。下表给出了这些试验的有关资料（把 $a$ 折算为 $1$）。

\begin{center}
\begin{tabular}{lccccc}
实验者 & 年份 & 针长 & 投掷次数 & 相交次数 & $\pi$ 的实验值 \\
\hline
Wolf & 1850 & 0.8 & 5\,000 & 2\,532 & 3.159\,6 \\
Smith & 1855 & 0.6 & 3\,204 & 1\,218.5 & 3.155\,4 \\
De Morgan, C. & 1860 & 1.0 & 600 & 382.5 & 3.137 \\
Fox & 1884 & 0.75 & 1\,030 & 489 & 3.159\,5 \\
Lazzerini & 1901 & 0.83 & 3\,408 & 1\,808 & 3.141\,592\,9 \\
Reina & 1925 & 0.541\,9 & 2\,520 & 859 & 3.179\,5
\end{tabular}
\end{center}

现在，随着计算机科学的发展，已按照上述思路建立起一类新的方法，称为\textbf{随机模拟法}，Buffon 投针实开其先河。
\end{example}

\subsection*{三、积分计算的 Monte Carlo 法}

二战后期，一批从事原子弹研究的美国科学家，与第一架数字电子计算机的诞生几乎同时提出了一系列随机模拟的计算方案。一大类是关于粒子运动的模拟，另一大类则是以积分计算为代表的新计算方法。他们还以赌城的名字命名这类方法，这就是著名的\textbf{Monte Carlo 方法}。下面介绍定积分计算的 Monte Carlo 法。

设在 $[a,b]$ 上，函数 $0\leqslant f(x)\leqslant M$，要求计算
\begin{align}
I=\int_a^b f(x)\,\mathrm{d}x\label{align:ch01-积分}
\end{align}
在图~\ref{figure:ch01-积分计算} 中，积分 $I$ 等于阴影所标的曲边梯形的面积。它被以 $\overline{ab}$ 为边，高为 $M$ 的矩形 $G$ 所包围，假如我们在 $G$ 中随机地取点，则该点落在阴影部分的概率为
\begin{align*}
P=\frac{I}{(b-a)M}
\end{align*}

\begin{figure}[H]
\centering
\includegraphics[width=0.5\textwidth]{ch01-montecarlo.png}
\caption{积分计算}\label{figure:ch01-积分计算}
\end{figure}

仿照 Buffon 投针问题的做法，投点 $N$ 次，计算落入阴影部分的次数 $n$，并以频率值 $\frac{n}{N}$ 作为概率 $P$ 之值代入上式，求得
\begin{align}
I=\frac{nM(b-a)}{N}\label{align:ch01-montecarlo}
\end{align}

这个做法不难推广到多重积分与任意边界的场合。

应当说明，当年设计的 Monte Carlo 新计算方法大都因效率不及传统方法的改良方案而遭淘汰，但是积分计算却因两大优点而生存：一是其误差与维数无关；二是适合于复杂的被积函数或边界。因此高维复杂积分的计算中，Monte Carlo 法仍有一席之地，而动态 Monte Carlo 方法则方兴未艾。Monte Carlo 方法的发展也大大推动了随机数的研究。

\subsection*{四、Bertrand 奇论}

几何概率在现代概率概念的发展中曾经起过重大作用。$19$ 世纪时，不少人相信，只要找到适当的等可能性描述，就可以给概率问题以唯一的解答，然而有人却构造出这样的例子，它包含着几种似乎都同样有理但却互相矛盾的答案，下面就是一个著名的例子。

\begin{example}[Bertrand 奇论]\label{example:ch01-bertrand奇论}
在半径为 $1$ 的圆内随机地取一条弦，问其长超过该圆内接等边三角形的边长 $\sqrt{3}$ 的概率等于多少？

这是一个几何概率问题，但是基于对术语“随机地”的含义的不同解释，这个问题却存在多种不同答案，下面是其中的三种。

\textbf{[解法一]}~~任何弦交圆周两点，不失一般性，先固定其中一点于圆周上，以此点为顶点作一等边三角形，显然只有落入此三角形内的弦才满足要求，这种弦的另一端跑过的弧长为整个圆周的 $\frac{1}{3}$，故所求概率等于 $\frac{1}{3}$（见图~\ref{figure:ch01-bertrand} (a)）。

\textbf{[解法二]}~~弦长只跟它与圆心的距离有关，而与方向无关，因此可以假定它垂直于某一直径。当且仅当它与圆心的距离小于 $\frac{1}{2}$ 时，其长才大于 $\sqrt{3}$，因此所求概率为 $\frac{1}{2}$（见图~\ref{figure:ch01-bertrand} (b)）。

\textbf{[解法三]}~~弦被其中点唯一确定，当且仅当其中点属于半径为 $\frac{1}{2}$ 的同心圆内时，弦长大于 $\sqrt{3}$，此小圆面积为大圆面积的 $\frac{1}{4}$，故所求概率等于 $\frac{1}{4}$（见图~\ref{figure:ch01-bertrand} (c)）。

\begin{figure}[H]
\centering
\includegraphics[width=0.9\textwidth]{ch01-bertrand.png}
\caption{Bertrand 奇论}\label{figure:ch01-bertrand}
\end{figure}

同一问题有三种不同的答案，细究其原因，发现是在取弦时采用不同的等可能性假定。在第一种解法中，假定端点在圆周上均匀分布，在第二种解法中则假定弦的中点在直径上均匀分布，而在第三种解法中又假定弦的中点在圆内均匀分布。这三种答案是针对三种不同的随机试验，对于各自的随机试验而言，它们都是正确的。

因此在使用术语“随机地”、“等可能”、“均匀分布”等时，应明确指明其含义；这又因试验而异。

$1899$ 年 Bertrand 在巴黎出版《概率论》，书中对几何概率提出了批评，并以这一生动的实例引起大家的注意。这种善意的批评，推动了概率论的发展。

由于采用等可能性来定义概率有这种困难，因此后来就选择另外的途径，即在定义概率这一基本概念时只指明概率应具有的基本性质，而把具体概率的给定放在一边。这样做的好处是能针对不同的随机试验给定适当的概率。

与概率的频率解释及古典概型一样，几何概率的研究对于我们了解应要求概率具有哪些基本性质是很有帮助的。
\end{example}

\subsection*{五、几何概率基本性质}

几何概率的定义及计算与几何图形的测度密切相关，因此所考虑的事件应是某种可定义测度的集合。这类集合的并、交也还应该是事件，甚至对它们的可列次并、交也应有这个要求。例如考察在 $(0,1)$ 中投一个点的随机试验，若以 $A$ 记该点落入 $\left(0,\frac{1}{2}\right)$ 中这个事件，而以 $A_n$ 记该点落入 $\left[\frac{1}{2^{n+1}},\frac{1}{2^n}\right)$ 中这一事件，$n=1,2,\cdots$ 则
\begin{align*}
A=\sum_{n=1}^{\infty}A_n
\end{align*}
如果假定所投的点落入某区间的概率等于该区间的长度，则 $P(A)=\frac{1}{2}$，而 $P(A_n)=\frac{1}{2^{n+1}}$，这时有
\begin{align*}
P(A)=\sum_{n=1}^{\infty}P(A_n)
\end{align*}
这里我们遇到了事件及概率的可列运算。

综上所述，几何概率应具有如下性质：

（i）对任何事件 $A$，$P(A)\geqslant 0$；

（ii）$P(\Omega)=1$；

（iii）若 $A_1,A_2,\cdots$ 两两互不相容，则
\begin{align}
P\left(\sum_{n=1}^{\infty}A_n\right)=\sum_{n=1}^{\infty}P(A_n)\label{align:ch01-可列可加}
\end{align}
前两个性质即非负性及规范性与古典型概率相同，\eqref{align:ch01-可列可加} 则要求对可列个两两互不相容事件成立，这性质称为\textbf{可列可加性}。

\end{document}